% !TEX root = userguide.tex

\def\headdate{1st June 2007}

\section{Convection-reaction equation using DG}

In this section we consider the convection-reaction equation using the
discontinuous Galerkin method (DG).
For an introduction to DG we refer to the 
book of Johnson \cite{Johnson87}.
We will only consider problems without diffusion. Diffusion can
be included with the Baumann-Oden approach
Baumann-Oden approach \cite{Baumann99}, but we will not do that here.

We consider only steady
state problems where all unknowns are solved at once, fully coupled. In a
time-dependent
approach this would correspond to a fully implicit approach.

\subsection{Weak form of the convection-reaction equation using discontinuous
 Galerkin \fem}

We consider the steady state convection-reaction equation in a region $\Omega$:
\begin{equation}
     \vec a\cdot\nabla u + b u = f\text{ on }\Omega
\end{equation}
where $\vek a$ is a divergence free velocity field ($\nabla\cdot\vek a=0$).
The boundary conditions are
\begin{gather}
    u = u_\text{in}\text{ on }\Gamma_\text{in}
\end{gather}
where the boundary $\Gamma_\text{in}$ is the inflow boundary ($\vek a\cdot\vek
n<0$).

For discontinuous Galerkin \fem\ we first start with a finite element mesh
before defining the weak form:
\[
\Omega=\cup_e\Omega_e, \quad\Omega_{e}\cap\Omega_{f}=\emptyset\quad
   \text{for }e\neq f
\]
The weak form can be now be stated, element wise, as:
find $u$ for all elements $e$, such that
\begin{equation}
 \int_{\Omega_e}v\Big(\vek a\nabla\cdot u+bu-f\Big)\D x
    +\int_{\Gamma^e_{\text{in}}}v\vek a\cdot\vek n[u]\D s=0,
        \quad\text{for all }v 
   \label{eq:DGweak}
\end{equation}
where $\Gamma^e_{\text{in}}$ is the inflow boundary of element $e$ and
$[u]=u^+-u$ is the ``jump'' across that element boundary. Also, $\vek
n$ is the outwardly directed unit normal vector on the boundary of
the element and $u^+$ is
the `outside' value of $u$.

We interpolate $u$ and $v$ by
polynomials on each element (approximation denoted by $u_h$ and $v_h$):
\begin{align*}
    u_h(\vek x,t)&=\tsum_i u_i(t)\phi_i(\vek x)=\col\phi^T(\vek x)\col u(t)
\\
    v_h(\vek x,t)&=\tsum_i v_i(t)\phi_i(\vek x)=\col\phi^T(\vek x)\col v(t)
\end{align*}
where $\col\phi(\vek x)$ are global shape functions. In contrast to Galerkin
\fem, the shape functions are \emph{discontinuous} across element boundaries.
This means that the first integral in Eq.~\eqref{eq:DGweak} only leads to
coupling
between unknowns \emph{within} an element. The boundary integral leads to a
coupling of element unknowns with neighbor elements, but in a non-standard
\fem\ way: all unknowns in an element are coupled with all unknowns of it's
neighbors.
 
Substituting the approximations into the weak form leads to
\[
 \col v^T\mat K\col u=
        \col v^T\col f \quad\text{for all }\col v
\]
or
\[
 \mat K\col u= \col f
\]
where the `stiffness matrix' $\mat K$ and
the right-hand side $\col f$ are given by the element matrices
\begin{align*}
 \mat K_e&=(\col\phi,\vek a\cdot\nabla{\col\phi^T}+b\col\phi)_{\Omega_e}
    -\int_{\Gamma^e_{\text{in}}}\vek a\cdot\vek n\col\phi\col\phi^T\D s
    +\int_{\Gamma^e_{\text{in}}}\vek a\cdot\vek n\col\phi\col\phi^{+T}\D s\\
\col f_e&= (\col\phi,f)_{\Omega_e}
\end{align*}
where $\phi^{+}$ are the shape functions in \emph{all} neighbor
elements\footnote{We only need to consider the boundary integration points
where $\vek a\cdot\vek n<0$, however \tfem\ assumes that there is a full
coupling between all neighbors (upstream and downstream). This has the
advantage that it is easy to program, diffusion can be included easily and the
resulting matrix has a symmetric pattern (preferable for the solvers). It
is however possible to limit the coupling to certain degrees of freedom, for
example only to the unknowns on the adjacent sides, greatly reducing the band
width of the matrix.} and
\[
   (a,b)_{\Omega_e}=\int_{\Omega_e}ab\D x
\]
Using index notation this becomes:
\begin{align*}
 K^e_{ij}&=(\phi_i,\vek a\cdot\nabla{\phi_j}+b\phi_j)_{\Omega_e}
    -\int_{\Gamma^e_{\text{in}}}\vek a\cdot\vek n\phi_i\phi_j\D s
    +\int_{\Gamma^e_{\text{in}}}\vek a\cdot\vek n\phi_i\phi^+_j\D s\\
f^e_i&= (\phi_i,f)_{\Omega_e}
\end{align*}
Note, that $\phi_j$ is only non-zero for $j$ in the current element $e$,
whereas $\phi^+_j$ is only non-zero for $j$ in the neighbor elements.
This means that the element matrix is \emph{non-square}. In \tfem\ the
assembling is split as follows:
\begin{enumerate}
   \item standard \fem\ assembly for the first term (volume integral) and the
         second term (boundary integral on the inside).
   \item assembly of the boundary integral on the outside (third term)
         is done for each side separately, i.e.\
         first side 1, then side 2, etc., each needing a different element
         matrix coupling the unknowns in the current element $e$ with a single
         adjacent element. 
\end{enumerate}
Note, that for an element on the boundary not all sides have an adjacent
element. If such a side is an inflow boundary, the value of $u^+=u_\text{in}$
must be specified as boundary condition and the contribution of the boundary
integral is to the right-hand side vector $\col f_e$ and not to the matrix.

We have opted for an element based assembly of the boundary integral. This
means that the integral is `scanned' twice, once from each element on a common
edge. Another option is to use an edge based description and integrate on all
edges and then compute the contribution to both elements on the common edge.
However, this would require a completely different bookkeeping that is
not available in \tfem.

\subsection{Example using an imposed exact solution}

We consider the steady state convection-reaction equation 
\begin{equation}
     \vec a\cdot\nabla u + u = f
\end{equation}
in a curved region with $\vec a=\vec e_1+2\vec e_2$. We impose the exact
solution
\begin{equation}
     u=1+\sin(\pi x)\sin(\pi y)
\end{equation}
by adjusting the right-hand side $f$ accordingly. The example can be obtained
by:
\begin{quote}
   \texttt{getexample convectionreaction}
\end{quote}
and the main program is \texttt{convectionreaction1.f90}. The main changes with
respect to the Poisson problem of Ch.~\ref{ch:examples} are:

\begin{listing}{47}
! problem definition
                                                                                
  call create_input_probdef ( mesh, input_probdef, nvec=1 )
                                                                                
  input_probdef%elementdof(1)%a = [(0,i=1,8),9]
  input_probdef%vec_elementdof(1)%a = reshape ( [(1,i=1,9)], [9,1] )
                                                                                
  call problem_definition ( input_probdef, mesh, problem )
                                                                                
\end{listing}
\begin{description}
   \item[line 51] all unknowns are located in the center nodal point (node 9)
since they are discontinuous across element boundaries.
   \item[line 52] the vector will be used for plotting the solution in a
continuous way.
\end{description}

\begin{listing}{61}
! create system matrix
                                                                                
  call create_sysmatrix_structure_base ( sysmatrix, mesh, problem )
                                                                                
  call create_sysmatrix_structure_dg ( sysmatrix, mesh, problem, &
    elemsubzeros=convecreac_sidelem_zeros )
                                                                                
  call finalize_sysmatrix_structure ( sysmatrix )
                                                                                
  call create_sysmatrix_data ( sysmatrix )
                                                                                
\end{listing}
\begin{description}
   \item[line 63--68] the structure of the matrix has been split in a standard
      \fem\ part, and a DG part.
   \item[line 63] here the structure of the standard \fem\ part is created, but
      it is not finalized. You cannot use the routine
    \texttt{create\_sysmatrix\_structure} here, because this includes the
     finalizing step.
   \item[line 65--66] here the DG structure is added. Note, that a parameter
    \texttt{elemsubzeros} is added, which supplies the name of an element
     subroutine that delivers an element matrix of type logical that says which
     components of the matrix are zero. In this particular
     case only the unknowns on the
     boundary need to be taken into account.
   \item[line 68] here the system matrix is finalized and no more additional
     structure can be added.
\end{description}

\begin{listing}{72}
! build system
                                                                                
  funcnr = 1
  funcnrinflow = 2
  vfuncnr = 1
                                                                                
  call build_system ( mesh, problem, sysmatrix, rhsd, elemsub=convecreac_elem )
                                                                                
  call build_system_dg ( mesh, problem, sysmatrix, elemsub=convecreac_sidelem,
&    addmatrix=.true., elemsubzeros=convecreac_sidelem_zeros )
                                                                                
  call check ( sysmatrix )
                                                                                
\end{listing}

\begin{description}
   \item[line 74-76] \texttt{funcnr} is for the right-hand side $f$,
     \texttt{funcnrinflow} is for the value of $u$ at the inflow boundary and
     \texttt{vfuncnr} is for the velocity $\vec a$.
   \item[line 78] the standard \fem\ assembly
   \item[line 80-81] the DG \fem\ assembly. Note that \texttt{addmatrix=.true.}
    has been added, avoiding setting everything to zero again. Note that also
    here the routine for the nonzero structure of the matrix has been added with
    \texttt{elemsubzeros}.
   \item[line 83] this is just a check that the matrix is indeed fully filled
      as specified in the structure of the matrix.
\end{description}

\begin{listing}{89}
! create exact solution
                                                                                
  call create_sysvector ( problem, solexact )
                                                                                
  funcnr = 2
                                                                                
  call fill_sysvector_elem ( mesh, problem, solexact, &
    elemsub=convecreac_fillexact )
                                                                                
\end{listing}

\begin{description}
   \item[line 95] this routine fills the sysvector element for element.
 This is 
 in a way similar to the assembly process for system vectors.
 However the difference is,
 that if a node appears twice the last appearance counts, since the data is just
 overwritten. Note, that an element routine must be written. Note also that
 this is the only way the exact solution can be filled in, because the
coordinate of the center node does not correspond to
 the position of the unknowns there and \text{fill\_sysvector} cannot be used.
\end{description}

\begin{listing}{102}
! compute vector for plotting

  call create_vector ( problem, vector, vec=1, elementwise=.true. )

  solution = sol

  call derive_vector ( mesh, problem, vector, elemsub=convecreac_vecelem )

! plot solution using figplot

  call plot_color_fill ( plot_options, mesh, problem, filename='sol_f.fig', &
    vector=vector )
  call plot_color_contour ( plot_options, mesh, problem, filename='sol_c.fig', &
    vector=vector )

\end{listing}
\begin{description}
\item[line 104] Here vector is created for plotting. The vector has been
defined elementwise (discontinuous across elements). Figplot can handle these
kind of vectors. For output to other postprocessing programs it may be
necessary to use a nodal averaged vector.
\end{description}

For details on the element implementation the reader is referred to the module 
\texttt{convecreac\_elements\_m} that is in the same directory as the main
program. An important structure for the implementation is
\texttt{mesh\%sidelem}, which contains information on the adjacent elements on
the sides of each element.

Run the problem by typing
\begin{quote}
   \texttt{convectionreaction1}
\end{quote}
and you will get something like
\begin{quote}
   \texttt{7.525107316978308E-03}
\end{quote}
on your screen. Typing the command
\begin{quote}
  \texttt{viewfig}
\end{quote}
will show something like Fig.~\ref{fig:vec}.
\begin{figure}[ht!]
   \begin{center}
   \includegraphics[scale=0.5]{vec}
   \end{center}
   \caption{Plot of $u$ for the convection-reaction equation}
    \label{fig:vec}
\end{figure}

\subsection{Example with periodic boundaries}

We use the same exact solution as in the previous example, however on a square
region $(0.2,2.2)\times(0.2,2.2)$. Furthermore, we use periodic boundaries
conditions in $x$ and $y$ direction. 
The main program is \texttt{convectionreaction2.f90}. The main change with
respect to the previous example is the mesh generation:

\begin{listing}{27}
! mesh generation
                                                                                
  meshgen_options%regionshape = 1
  meshgen_options%lx = 2.0_dp
  meshgen_options%ly = 2.0_dp
  meshgen_options%ox = 0.3_dp
  meshgen_options%oy = 0.3_dp
                                                                                
  nx = 20
  ny = 20
                                                                                
  meshgen_options%elshape = 6
  meshgen_options%nx = nx
  meshgen_options%ny = ny
                                                                                
  call quadrilateral2d ( mesh, meshgen_options )
                                                                                
  call add_to_mesh ( mesh, curve=[-3] )  ! curve 5
  call add_to_mesh ( mesh, curve=[-4] )  ! curve 6
                                                                                
  call glue_mesh ( mesh, curve1=1, curve2=5 )
  call glue_mesh ( mesh, curve1=2, curve2=6 )
                                                                                
  call fill_mesh_parts ( mesh )

\end{listing}

\begin{description}
   \item[lines 29--42] generation of the square mesh
   \item[lines 44--45] two new curves (C5 and C6) are added here, which are
      defined in opposite direction to curves C3 and C4 (see
      Fig.~\ref{fig:unitsquare}), respectively.
   \item[line 47] curves C1 and C5 are glued together. This means that for
     \dgfem\ these curves are the same: an element having a side on C1 is
     assumed to have the corresponding element on C5 as its neighbor and vice
     verse. So if there is inflow on curve C1 it comes from the outflow of C5
     as if the periodic regions are all stacked to fill the infinite domain.
    The two curves that are connected\footnote{Note that 'connected' has
    a different meaning here than the connection using \texttt{meshconnect} in
    \sep. In \sep\ the nodes are connected such that they have the same 
    degrees of freedom. Here the routine \texttt{glue\_mesh} will only change
    the `neighbors' of the elements, i.e.\ it changes the structure
\texttt{mesh\%sidelem}, which contains information on the adjacent elements on
the sides of each element.
    Therefore the term 'glue' has been chosen to
    emphasize the distinction.} should be in the same direction.
   \item[line 48] curves C2 and C6 are glued together 
   \item[line 50] the routine \texttt{fill\_mesh\_parts} must be called last.
\end{description}

Run the problem by typing
\begin{quote}
   \texttt{convectionreaction2}
\end{quote}
and you will get something like
\begin{quote}
   \texttt{6.052686362307824E-04}
\end{quote}
on your screen. Typing the command
\begin{quote}
  \texttt{viewfig}
\end{quote}
will show something like Fig.~\ref{fig:vec2}.
\begin{figure}[ht!]
   \begin{center}
   \includegraphics[scale=0.5]{vec2}
   \end{center}
   \caption{Plot of $u$ for the convection-reaction equation with periodical
            boundary conditions.}
    \label{fig:vec2}
\end{figure}

